% ==============================================================================
% AGRADECIMENTOS (opcional; manual, cap. 2, item 7)
% - Título "AGRADECIMENTOS" centralizado e em negrito (gerado pela classe).
% - Vem após a dedicatória (ou após a folha de aprovação, se não houver
%   dedicatória), em folha própria, sem número impresso.
% - Texto em parágrafos normais, com recuo. Não usar \noindent nem \newline
%   para separar parágrafos: basta deixar uma linha em branco entre eles.
% - Ordem sugerida: orientação, banca, docentes, colegas, instituição e
%   agências de fomento, família e amigos. A ordem é livre, mas convém
%   manter as menções institucionais juntas.
% - Tom pessoal, porém sóbrio. Evitar exageros, piadas internas e listas
%   extensas de nomes. Uma página costuma ser suficiente.
% - Usar o tratamento correto de cada pessoa: Prof. Dr. / Profa. Dra.
% - Apagar os parágrafos que não se aplicam e todos os campos \preencher.
% ==============================================================================
\begin{agradecimentos}

  % --- Orientação ---------------------------------------------------------------
  % Mencionar contribuições concretas (rigor, disponibilidade, autonomia),
  % e não apenas elogios genéricos.
  Ao meu orientador, \preencher{Prof. Dr. nome completo}, pela orientação
  criteriosa, pela confiança e pelas contribuições decisivas ao desenvolvimento
  deste trabalho.

  % Excluir se não houver coorientação.
  Ao meu coorientador, \preencher{Prof. Dr. nome completo}, pelas sugestões
  e pelo acompanhamento ao longo da pesquisa.

  % --- Banca e docentes ---------------------------------------------------------
  % Na versão de Defesa, agradecer aos membros da banca de qualificação, cujas
  % recomendações foram incorporadas ao texto.
  Aos professores \preencher{nomes}, membros da banca de qualificação, pelas
  recomendações que aprimoraram esta \preencher{dissertação ou tese}.

  % Pode ser geral ou nominal; se nominal, citar poucos nomes e justificar.
  Aos docentes do Programa de Pós-Graduação em \preencher{Informática e Gestão
    do Conhecimento ou Gestão de Projetos}, pelo conhecimento compartilhado
  nas disciplinas e seminários.

  % --- Colegas ------------------------------------------------------------------
  Aos colegas do \preencher{grupo de pesquisa ou turma}, pela troca de ideias
  e pelo apoio mútuo.

  % --- Instituição e fomento ----------------------------------------------------
  % Bolsistas FAP/UNINOVE:
  Agradeço à Universidade Nove de Julho (UNINOVE) pela concessão da bolsa de
  estudos do Fundo de Apoio à Pesquisa (FAP).

  % Outras agências: atenção à crase (à CAPES, à FAPESP, mas ao CNPq) e ao
  % número do processo, exigido pelo CNPq e pelas fundações estaduais.
  Agradeço \preencher{à CAPES, ao CNPq ou à agência de fomento} pela
  concessão \preencher{da bolsa de estudos ou da taxa}
  (processo n.\textsuperscript{o}~\preencher{número}).

  % Bolsistas CAPES: frase OBRIGATÓRIA (Portaria CAPES n.º 206/2018), com
  % esta redação exata. Descomentar se aplicável.
  % O presente trabalho foi realizado com apoio da Coordenação de
  % Aperfeiçoamento de Pessoal de Nível Superior -- Brasil (CAPES) --
  % Código de Financiamento 001.

  % --- Família e amigos --------------------------------------------------------
  À minha família, especialmente \preencher{nomes}, pelo apoio incondicional
  e pela compreensão nos momentos de ausência.

  % --- Fechamento ---------------------------------------------------------------
  Por fim, agradeço a todos que, direta ou indiretamente, contribuíram para a
  realização deste trabalho.

\end{agradecimentos}